\textbf{Phần 1. Trắc nghiệm lựa chọn}
\setcounterpageref{ex}{0}

\begin{ex} %[2D2V09]
Kết quả của phép tính $15 - 8$ là
\choice
{\True $7$}
{$8$}
{$6$}
{$9$}
\loigiai{
$15 - 8 = 7$.
\textbf{Đáp án của câu hỏi là A}
}
\end{ex}

\begin{ex} %[2D2V09]
Một cành cây có $14$ con chim đậu. Bắn một tiếng súng làm $6$ con chim bay đi. Lát sau lại có $3$ con chim bay đi nữa. Hỏi trên cành cây còn lại bao nhiêu con chim?
\choice
{$6$ con chim}
{$4$ con chim}
{\True $5$ con chim}
{$7$ con chim}
\loigiai{
Số chim còn lại: $14 - 6 - 3 = 8 - 3 = 5$ (con chim).
\textbf{Đáp án của câu hỏi là C}
}
\end{ex}

\begin{ex} %[2D2V09]
Tìm $x$ biết: $13 - x = 7$.
\choice
{$5$}
{\True $6$}
{$7$}
{$8$}
\loigiai{
$x = 13 - 7 = 6$.
\textbf{Đáp án của câu hỏi là B}
}
\end{ex}

\begin{ex} %[2D2V09]
Bà có $17$ quả trứng gà. Bà đã dùng $9$ quả để làm bánh. Số trứng gà bà còn lại là
\choice
{$9$ quả}
{$7$ quả}
{$6$ quả}
{\True $8$ quả}
\loigiai{
Số trứng gà còn lại: $17 - 9 = 8$ (quả).
\textbf{Đáp án của câu hỏi là D}
}
\end{ex}

\begin{ex} %[2D2V09]
Tính giá trị của biểu thức: $16 - 9 + 5$.
\choice
{\True $12$}
{$11$}
{$13$}
{$10$}
\loigiai{
$16 - 9 = 7$; $7 + 5 = 12$.
\textbf{Đáp án của câu hỏi là A}
}
\end{ex}

\begin{ex} %[2D2V09]
Mẹ mua $12$ bông hoa hồng và hoa cúc, trong đó có $7$ bông hoa hồng. Hỏi mẹ mua bao nhiêu bông hoa cúc?
\choice
{$6$ bông hoa}
{$4$ bông hoa}
{\True $5$ bông hoa}
{$7$ bông hoa}
\loigiai{
Số bông hoa cúc: $12 - 7 = 5$ (bông hoa).
\textbf{Đáp án của câu hỏi là C}
}
\end{ex}

\begin{ex} %[2D2V09]
Điền dấu thích hợp: $14 - 7 \dots 13 - 6$.
\choice
{$>$}
{\True $=$}
{$<$}
{Không so sánh được}
\loigiai{
$14 - 7 = 7$; $13 - 6 = 7$. Hai vế bằng nhau nên điền dấu $=$.
\textbf{Đáp án của câu hỏi là B}
}
\end{ex}

\begin{ex} %[2D2V09]
Hiệu giữa số chẵn lớn nhất trong phạm vi $20$ và số lẻ lớn nhất có 1 chữ số là
\choice
{$10$}
{$12$}
{$9$}
{\True $11$}
\loigiai{
Số chẵn lớn nhất trong phạm vi 20 là 20. Số lẻ lớn nhất có 1 chữ số là 9.\nHiệu: $20 - 9 = 11$.
\textbf{Đáp án của câu hỏi là D}
}
\end{ex}

\begin{ex} %[2D2V09]
Tìm một số biết rằng lấy $18$ trừ đi số đó rồi cộng thêm $4$ thì được $13$.
\choice
{\True $9$}
{$8$}
{$10$}
{$7$}
\loigiai{
Ta có: $18 - x = 13 - 4 = 9 \Rightarrow x = 18 - 9 = 9$.
\textbf{Đáp án của câu hỏi là A}
}
\end{ex}

\begin{ex} %[2D2V09]
Một bác thợ làm được $16$ sản phẩm trong hai ngày. Ngày thứ nhất bác làm được $9$ sản phẩm. Hỏi ngày thứ nhất bác làm nhiều hơn ngày thứ hai bao nhiêu sản phẩm?
\choice
{$7$ sản phẩm}
{$3$ sản phẩm}
{\True $2$ sản phẩm}
{$1$ sản phẩm}
\loigiai{
Ngày thứ hai bác làm: $16 - 9 = 7$ (sản phẩm).\nNgày thứ nhất hơn ngày thứ hai: $9 - 7 = 2$ (sản phẩm).
\textbf{Đáp án của câu hỏi là C}
}
\end{ex}

\begin{ex} %[2D2V09]
Kết quả của phép trừ $11 - 4$ bằng kết quả của phép tính nào dưới đây?
\choice
{$15 - 7$}
{\True $16 - 9$}
{$13 - 5$}
{$12 - 6$}
\loigiai{
$11 - 4 = 7$. Ta thấy $16 - 9 = 7$. Do đó chọn $16 - 9$.
\textbf{Đáp án của câu hỏi là B}
}
\end{ex}

\begin{ex} %[2D2V09]
Trong rổ có $15$ quả cam. Mai lấy ra $6$ quả, Nam lấy ra $4$ quả. Hỏi trong rổ còn lại bao nhiêu quả cam?
\choice
{$6$ quả}
{$7$ quả}
{$4$ quả}
{\True $5$ quả}
\loigiai{
Hai bạn lấy ra: $6 + 4 = 10$ (quả).\nTrong rổ còn: $15 - 10 = 5$ (quả).
\textbf{Đáp án của câu hỏi là D}
}
\end{ex}

\textbf{Phần 2. Trắc nghiệm đúng sai}
\setcounterpageref{ex}{0}

\begin{ex} %[2D2V09]
Xét tính đúng sai của các phép trừ có nhớ trong phạm vi 20:
\choiceTF
{\True Phép trừ $14 - 8$ có hiệu bằng $6$.}
{Hiệu của $15 - 7$ bằng hiệu của $16 - 9$.}
{\True Trong phép trừ $17 - 9 = 8$, số bị trừ là 17, số trừ là 9.}
{\True Nếu lấy 18 trừ đi một số có một chữ số thì hiệu luôn lớn hơn hoặc bằng 9.}
\loigiai{
\begin{itemize}
    \item Ý a) Đúng: $14 - 8 = 6$.
    \item Ý b) Sai: $15 - 7 = 8$, còn $16 - 9 = 7$. Hai kết quả khác nhau.
    \item Ý c) Đúng: Đúng theo tên gọi các thành phần.
    \item Ý d) Đúng: Số trừ lớn nhất là 9 nên $18 - 9 = 9$. Hiệu luôn $\ge 9$.
\end{itemize}
\begin{itemize}
    \item a) \textbf{Đúng}.
    \item b) \textbf{Sai}.
    \item c) \textbf{Đúng}.
    \item d) \textbf{Đúng}.
\end{itemize}
}
\end{ex}

\begin{ex} %[2D2V09]
Một nhóm bạn gồm 15 bạn chơi trò trốn tìm, trong đó đã có 8 bạn bị tìm thấy:
\choiceTF
{\True Còn lại 7 bạn chưa bị tìm thấy.}
{Số bạn đã bị tìm thấy ít hơn số bạn chưa bị tìm thấy 1 bạn.}
{\True Nếu tìm thấy thêm 3 bạn nữa thì số bạn chưa bị tìm thấy là 4 bạn.}
{Số bạn bị tìm thấy gấp đôi số bạn chưa bị tìm thấy.}
\loigiai{
\begin{itemize}
    \item Ý a) Đúng: $15 - 8 = 7$ bạn.
    \item Ý b) Sai: 8 bạn bị tìm thấy nhiều hơn 7 bạn chưa tìm thấy 1 bạn.
    \item Ý c) Đúng: $7 - 3 = 4$ bạn.
    \item Ý d) Sai: $8$ không gấp đôi $7$. Khẳng định sai.
\end{itemize}
\begin{itemize}
    \item a) \textbf{Đúng}.
    \item b) \textbf{Sai}.
    \item c) \textbf{Đúng}.
    \item d) \textbf{Sai}.
\end{itemize}
}
\end{ex}

\textbf{Phần 3. Trả lời ngắn}
\setcounterpageref{ex}{0}

\begin{ex} %[2D2V09]
Tính giá trị của biểu thức: $17 - 8 + 6$.
\shortans{15}
\loigiai{
\textbf{Đáp số:} $15$\\
$17 - 8 = 9$; $9 + 6 = 15$.
}
\end{ex}

\begin{ex} %[2D2V09]
Trong bể có 16 con cá. Bố vớt ra 7 con cá đem tặng bạn. Hỏi trong bể còn lại bao nhiêu con cá?
\shortans{9}
\loigiai{
\textbf{Đáp số:} $9$\\
Số cá còn lại: $16 - 7 = 9$ con.
}
\end{ex}

\begin{ex} %[2D2V09]
Tìm số tự nhiên x biết: $15 - x = 8$.
\shortans{7}
\loigiai{
\textbf{Đáp số:} $7$\\
$x = 15 - 8 = 7$.
}
\end{ex}

\begin{ex} %[2D2V09]
Có hai ngăn kéo đựng tất cả 18 quyển sách. Ngăn thứ nhất có 9 quyển sách. Hỏi ngăn thứ nhất có nhiều hơn hay ít hơn ngăn thứ hai bao nhiêu quyển sách?
\shortans{0}
\loigiai{
\textbf{Đáp số:} $0$\\
Ngăn 2 có: $18 - 9 = 9$ quyển. Hai ngăn bằng nhau nên hiệu bằng 0.
}
\end{ex}

\begin{ex} %[2D2V09]
Một cây tre dài 14 m bị bão làm gãy một đoạn ngọn dài 6 m. Hỏi đoạn gốc còn lại dài bao nhiêu mét?
\shortans{8}
\loigiai{
\textbf{Đáp số:} $8$\\
Đoạn gốc dài: $14 - 6 = 8\text{ m}$.
}
\end{ex}

\begin{ex} %[2D2V09]
Tìm hiệu của số tròn chục nhỏ nhất có hai chữ số và số lẻ nhỏ nhất có 1 chữ số khác 1.
\shortans{7}
\loigiai{
\textbf{Đáp số:} $7$\\
Số tròn chục nhỏ nhất là 10. Số lẻ nhỏ nhất khác 1 là 3. Hiệu là: $10 - 3 = 7$.
}
\end{ex}
